% !TEX root = ../main.tex
\section{Logical coherence and decision theory}

\red{Connect the choice of $\phi$ to decision theory.}

Does there exist a utility function such that,
if $\phi^* = \arg \max_{\phi} \E[U(\phi,\theta)]$,
then $\phi^*$ is $\aleph_0$-coherent?

\begin{definition}
 $P_{\phi}(x) \subset \Theta$ and
 $N_{\phi}(x) \subset \Theta$ are such that
 \label{def:region_p_n}
 \begin{align*}
  P_{\phi}(x) 
  &= \left\{\theta \in \Theta:
  \forall H \in \sigma(\Theta),
  \theta \in H \rightarrow 
  \phi(H|x) \geq \half\right\}, \text{ and } \\
  N_{\phi}(x)
  &= \left\{\theta \in \Theta:
  \forall H \in \sigma(\Theta),
  \theta \notin H \rightarrow 
  \phi(H|x) \leq \half\right\}.
 \end{align*}
\end{definition}

$P_{\phi}$ and $N_{\phi}$ can be
interpreted as two possible definitions of
sets of values of $\theta$ that are
highly plausible according to $\phi$.
If $\theta_0 \in P_{\phi}$, then
$\theta_0 \in H$ is a sufficient condition
for $H$ not to be taken as false.
That is, $\theta_0$ is so plausible that
its presence in $H$ is sufficient
to avoid $H$ being taken as false.
Similarly, if $\theta_0 \in N_{\phi}$, then
$\theta_0 \in H$ is a necessary condition for
$H$ to be taken as true.
That is, $\theta_0$ is so plausible that,
if $\theta \notin H$, then
one does not take $H$ to be true.

\begin{definition}
 \label{def:u_gfbst}
  Let $\sP(\sigma(\Theta))$ be
  the powerset of $\sigma(\Theta)$ and
  $\rho: \mathcal{P}(\sigma(\Theta))
 \rightarrow \leftindex^*{\Re}$ be
 a nonstandard probability over
 collections of hypotheses such that,
 if $H$ is a proper subset of $H'$, then
 $\rho(H) < \rho(H')$.
 The existence of $\mu$ follows from
 \citet{Schervish2022}. Also, let
 $\rho$ be a measure over $\sigma(\Theta)$.
 \begin{align*}
  U(\phi, \theta)
  &:= \I(\theta \in P_{\phi}(X))
 + \I(\theta \in N_{\phi}(X))
  -\lambda \left(\mu(P_{\phi}(X))
  +\mu(N_{\phi}(X)) \right) 
  -\odot \cdot \rho(A_{\phi}(X)), \\
  & \text{where } A_{\phi}(x) = 
  \left\{H \in \sigma(\Theta): \phi(H|x) = \half 
  \right\} \text{and $\odot$ is an infinitesimal}.
 \end{align*}
\end{definition}

\Cref{def:u_gfbst} admits 
an intuitive interpretation.
Recall that $P_{\phi}$ and $N_{\phi}$ are
sets of plausible values for $\theta$.
The first two summands,
$\I(\theta \in P_{\phi}(X)) + 
\I(\theta \in N_{\phi}(X))$, reward
when $\theta \in P_{\phi}$ and
when $\theta \in P_{\phi}$. Also,
$-\lambda \left(\mu(P_{\phi}(X))
  +\mu(N_{\phi}(X)) \right)$ punishes
when $P_{\phi}$ and $N_{\phi}$ are large sets.
Finally, $-\odot \cdot \rho(U_{\phi})$ puts 
a small penalty for each hypothesis
in which $\phi$ remains undecided.

\begin{theorem}
 \label{thm:min_u}
 $\phi^* = \arg \max_{\phi} \E[U(\phi,\theta)]$
 if and only if
 $\phi^*$ is a version of $\phi_0$, where
 $\phi_0$ is based on the region, $R_x$,
 and $R_x$ is
 the highest density set:
 \begin{align*}
  \left\{\theta \in \Theta: 
  \frac{d\P(\theta|x)}{d \mu} 
  > (1+\lambda)^{-1}\right\}.
 \end{align*}
\end{theorem}
